% !PRINTED 551
\weberpage{551}
If, as in \S.\,163, we represent each root in the sequence of roots
$\alpha, \alpha_1, \ldots \alpha_{m-1}$ rationally by the preceding one:
\[
\alpha_1 = \Theta(\alpha), \quad \alpha_2 = \Theta(\alpha_1), \ldots \alpha_{m-1} = \Theta(\alpha_{m-2}), \quad \alpha = \Theta(\alpha_{m-1}),
\]
in which, since $\Omega$ is here assumed to be real, $\Theta(x)$ denotes a real rational function of $x$, then it follows first that if
one of the roots is real, all the others must be real as well.
This is always the case for odd $m$, since a real equation of odd degree must always have at least one real
root.

For even $m$ imaginary roots may also be present,
and if one root is imaginary, then all must be so,
since, if a real root occurs, all the others are also
real. If we denote by $\Theta^{\nu}(x)$ the $\nu$-fold repetition of the
function $\Theta$, then
% [sic] authority prints alpha^k; the later parallel occurrence alpha_k confirms k should be a subscript.
\[
\alpha_{\nu+k} = \Theta^{\nu}(\alpha^k).
\]
If then $\alpha = \alpha_0$ is conjugate imaginary with $\alpha_k$, then for
every $\nu$ the functions $\Theta^{\nu}(\alpha_0)$ and $\Theta^{\nu}(\alpha_k)$, that is, $\alpha_{\nu}$ and $\alpha_{\nu+k}$,
are also conjugate imaginary.

From this it follows that $2k = m$ must hold, and we conclude
that in the case of imaginary roots
\[
\alpha_k \quad \text{and} \quad \alpha_{k + \frac{m}{2}}
\]
form, for every index $k$, a pair of conjugate imaginary roots.

The cubic equations become cyclic equations by the adjunction of the
square root of the discriminant.
When the discriminant is positive, the roots are, in
agreement with this theorem, real.

\begin{center}
\S.\,166.\\
\gesperrt{Division of the angle.}
\end{center}

Cyclic equations also include those
on which the division of an angle into $m$ equal parts depends; to these we have reduced the general cyclic equations in
a real field.

The problem can be formulated thus:

\gesperrt{When $\cos m\varphi$ and $\sin m\varphi$ are given, then
$\cos\varphi$ and $\sin\varphi$ are to be found from them.}
